\documentclass[12pt, a4paper]{article}
\usepackage[utf8]{inputenc}
\usepackage[margin=0.9in]{geometry}
\usepackage{amsmath, amssymb, amsfonts}
\usepackage{xcolor}
\usepackage{listings}
\usepackage{tikz}
\usepackage{titlesec}
\usepackage{hyperref}
\usepackage{tcolorbox}
\usepackage{booktabs}
\usepackage{enumitem}
\usepackage{pagecolor}
\usepackage{fancyhdr}
\usepackage{setspace}
\usepackage{charter} % Elegant font
\usepackage{tocloft}
\usepackage{graphicx}
\usepackage{fontawesome5} % For icons
\usepackage{multicol}
\usepackage[most]{tcolorbox}
% --- COMPLETELY UNIQUE COLOR PALETTE ---
\definecolor{DeepForest}{RGB}{20, 60, 40}      % Rich forest green
\definecolor{AmberSunset}{RGB}{220, 110, 40}   % Warm amber-orange
\definecolor{DarkCharcoal}{RGB}{45, 45, 50}    % Deep gray for text
\definecolor{Ivory}{RGB}{255, 252, 245}        % Warm off-white
\definecolor{MistyLavender}{RGB}{230, 225, 240} % Soft lavender for code
\definecolor{CrimsonRed}{RGB}{190, 50, 50}
\definecolor{TealBlue}{RGB}{30, 120, 130}

% --- Apply Unique Background ---
\pagecolor{Ivory}

% --- Global Styling ---
\renewcommand{\familydefault}{\rmdefault} 
\hypersetup{colorlinks=true, linkcolor=DeepForest, urlcolor=AmberSunset, citecolor=AmberSunset}

% --- MODERN GEOMETRIC SECTION HEADERS ---
\titleformat{\section}
    {\color{DeepForest}\Large\bfseries\sffamily}
    {\makebox[1.5cm][l]{\color{AmberSunset}\thesection}}{0em}
    {}
    [\vspace{-0.2cm}\rule{\textwidth}{0.3pt}]

\titleformat{\subsection}
    {\color{DeepForest}\large\bfseries\sffamily}
    {\color{AmberSunset}\thesubsection}{0.5em}{}
    
\titleformat{\subsubsection}
    {\color{DarkCharcoal}\normalsize\bfseries\itshape}
    {\color{TealBlue}\thesubsubsection}{0.5em}{}

% --- CREATIVE HEADER/FOOTER ---
\pagestyle{fancy}
\fancyhf{}
\fancyhead[L]{\color{DeepForest}\faChartLine \quad \textit{Statistical Data Analysis Report}}
\fancyhead[R]{\color{DeepForest}\faFileAlt \quad Page \thepage}

\renewcommand{\headrulewidth}{0.8pt}
\renewcommand{\headrule}{\color{AmberSunset!60}\hrule width\headwidth height\headrulewidth}

% --- MINIMALIST TOC ---
\renewcommand{\cftsecleader}{\cftdotfill{\cftdotsep}}
\renewcommand{\cftsecfont}{\color{DeepForest}\bfseries\sffamily}
\renewcommand{\cftsecpagefont}{\color{DeepForest}}
\renewcommand{\cftsubsecfont}{\color{DarkCharcoal}}
\setlength{\cftbeforesecskip}{5pt}

% --- UNIQUE CODE BLOCK - MINIMALIST WITH LEFT BAR ---
\lstset{
    language=Python,
    basicstyle=\ttfamily\small\color{DarkCharcoal},
    breaklines=true,
    backgroundcolor=\color{MistyLavender!30},
    frame=none,
    keywordstyle=\color{DeepForest}\bfseries,
    commentstyle=\color{TealBlue}\itshape,
    stringstyle=\color{CrimsonRed},
    xleftmargin=20pt,
    showstringspaces=false,
    numbers=left,
    numberstyle=\tiny\color{gray},
    numbersep=10pt,
    literate={\ \ }{{\ }}1,
    rulecolor=\color{AmberSunset!40}
}

% Add left bar to code blocks
\BeforeBeginEnvironment{lstlisting}{\begin{tcolorbox}[colback=MistyLavender!30, colframe=AmberSunset!60, boxrule=0pt, leftrule=3pt, arc=0pt, enhanced, outer arc=0pt]}
\AfterEndEnvironment{lstlisting}{\end{tcolorbox}}

% --- CREATIVE OUTPUT BOX ---
\newtcolorbox{outputbox}{
    colback=white,
    colframe=TealBlue!50,
    boxrule=0.5pt,
    arc=8pt,
    left=12pt,
    right=12pt,
    top=10pt,
    bottom=10pt,
    fontupper=\small\ttfamily,
    before upper={\color{DeepForest}\faTerminal\quad \textbf{RESULT:}\vspace{0.2cm}\par},
    enhanced,
    drop fuzzy shadow=AmberSunset!30,
    borderline={0.5pt}{0pt}{AmberSunset!40}
}

% --- BOLD PROBLEM HEADER - CARD STYLE ---
\newcounter{problemcounter}
\setcounter{problemcounter}{1}

\newcommand{\problemtitle}[1]{%
    \vspace{1cm}
    \begin{tcolorbox}[
        colback=DeepForest!5,
        colframe=DeepForest,
        arc=15pt,
        boxrule=1.2pt,
        left=20pt,
        right=20pt,
        top=12pt,
        bottom=12pt,
        enhanced,
        attach boxed title to top center={yshift=-8pt},
        boxed title style={
            colback=AmberSunset,
            colframe=DeepForest,
            arc=10pt,
            outer arc=10pt,
            boxrule=0.8pt,
            sharp corners,
            },
        title={\color{white}\large\sffamily\bfseries \faTasks \quad PROBLEM \arabic{problemcounter}}
    ]
        \centering
        {\color{DeepForest}\Large\bfseries\sffamily #1}
    \end{tcolorbox}
    \stepcounter{problemcounter}
    \vspace{0.3cm}
}

% --- KEY INSIGHT BOX ---
\newtcolorbox{insightbox}{
    colback=TealBlue!8,
    colframe=TealBlue!60,
    arc=8pt,
    boxrule=0.8pt,
    left=12pt,
    right=12pt,
    top=10pt,
    bottom=10pt,
    fontupper=\normalsize,
    before upper={\color{TealBlue}\faLightbulb\quad \textbf{Key Insight}\par\medskip},
    enhanced,
    drop fuzzy shadow
}

% --- STATISTICAL NOTE BOX ---
\newtcolorbox{statbox}{
    colback=AmberSunset!5,
    colframe=AmberSunset!50,
    arc=5pt,
    boxrule=0.5pt,
    left=10pt,
    right=10pt,
    top=8pt,
    bottom=8pt,
    fontupper=\small\itshape,
    before upper={\color{AmberSunset}\faInfoCircle\quad },
}

% --- SIDEBAR FOR IMPORTANT NUMBERS ---
\newtcolorbox{numberbox}[1]{
    colback=DeepForest!85,
    coltext=white,
    arc=8pt,
    boxrule=0pt,
    left=8pt,
    right=8pt,
    top=6pt,
    bottom=6pt,
    fontupper=\bfseries\large,
    width=2cm,
    halign=center,
    before upper={\faHashtag\quad #1\\},
}

\begin{document}

%% ===================== MINIMALIST COVER PAGE =====================
\pagestyle{empty}

\begin{tikzpicture}[remember picture, overlay]
    % Clean white background with subtle texture
    \fill[Ivory] (current page.south west) rectangle (current page.north east);
    
    % Diagonal accent stripe from top-left
    \fill[DeepForest!15] (current page.north west) -- ([xshift=5cm]current page.north west) -- 
                         ([yshift=-8cm]current page.north east) -- ([xshift=-5cm, yshift=-8cm]current page.north east) -- cycle;
    
    % Bottom-right accent
    \fill[AmberSunset!12] (current page.south east) -- ([xshift=-6cm]current page.south east) -- 
                          ([yshift=7cm]current page.south west) -- ([xshift=6cm, yshift=7cm]current page.south west) -- cycle;
    
    % Large geometric circles
    \draw[AmberSunset!20, line width=3pt] ([xshift=-2cm, yshift=-3cm]current page.center) circle (5cm);
    \draw[DeepForest!15, line width=2pt] ([xshift=3cm, yshift=4cm]current page.center) circle (3.5cm);
    
    % University Name - Clean and centered
    \node[anchor=north] at ([yshift=-2.5cm]current page.north) {
        {\color{DeepForest}\fontsize{20}{24}\selectfont \sffamily \textbf{BEGUM ROKEYA UNIVERSITY, RANGPUR}}
    };
    
    % Decorative line
    \node[anchor=north] at ([yshift=-3.5cm]current page.north) {
        {\color{AmberSunset}\rule{4cm}{1.5pt}}
    };
    
    \node[anchor=north] at ([yshift=-4cm]current page.north) {
        {\color{DarkCharcoal}\large \textsc{Department of Statistics}}
    };
    
    % Main Title - Bold and modern
    \node[anchor=center] at ([yshift=1.5cm]current page.center) {
        \begin{minipage}{13cm}
            \centering
            {\color{DeepForest}\fontsize{38}{44}\selectfont \sffamily \textbf{STATISTICAL}}\\[0.2cm]
            {\color{DeepForest}\fontsize{38}{44}\selectfont \sffamily \textbf{DATA ANALYSIS}}\\[0.5cm]
            {\color{AmberSunset}\fontsize{16}{20}\selectfont \textit{A Comprehensive Laboratory Report}}\\[0.8cm]
        \end{minipage}
    };
    
    % === COURSE DETAILS - MOVED TO MIDDLE ===
    \node[anchor=center] at ([yshift=-2.5cm]current page.center) {
        \begin{tcolorbox}[
            colback=DeepForest!8,
            colframe=AmberSunset!70,
            arc=12pt,
            boxrule=1.2pt,
            left=15pt,
            right=15pt,
            top=10pt,
            bottom=10pt,
            enhanced,
            drop fuzzy shadow=AmberSunset!30,
            width=10cm
        ]
            \centering
            {\color{DeepForest}\large\sffamily\bfseries \faBook \quad COURSE DETAILS}\\[0.3cm]
            {\color{AmberSunset}\rule{6cm}{0.5pt}}\\[0.3cm]
            {\color{DeepForest}\Large\bfseries\sffamily STAT-4104}\\[0.2cm]
            {\color{DarkCharcoal}\large \textsc{Categorical Data Analysis}}
        \end{tcolorbox}
    };
    
    % Author Info - Minimalist table
    \node[anchor=south] at ([yshift=3cm]current page.south) {
        \begin{minipage}{12cm}
            \centering
            \renewcommand{\arraystretch}{1.3}
            \begin{tabular}{r l | r l}
                \multicolumn{2}{c}{\textbf{\color{DeepForest}\sffamily SUBMITTED BY}} & \multicolumn{2}{c}{\textbf{\color{DeepForest}\sffamily SUBMITTED TO}} \\[0.2cm]
                \textcolor{DarkCharcoal}{Name:} & \textcolor{DeepForest}{Prome Saha} & \textcolor{DarkCharcoal}{Name:} & \textcolor{DeepForest}{Dr. Md. Siddikur Rahman} \\
                \textcolor{DarkCharcoal}{ID:} & 12110042 & \textcolor{DarkCharcoal}{Title:} & Associate Professor \\
                \textcolor{DarkCharcoal}{Reg No:} & 000015905 & \textcolor{DarkCharcoal}{Dept:} & Statistics \\
                \textcolor{DarkCharcoal}{Session:} & 2021-2022 & \textcolor{DarkCharcoal}{Inst:} & Begum Rokeya University \\
                \textcolor{DarkCharcoal}{Year/Sem:} & 4\textsuperscript{th} Year, 1\textsuperscript{st} Sem & & 
            \end{tabular}
        \end{minipage}
    };
    
    % Date - Bottom
    \node[anchor=south] at ([yshift=1cm]current page.south) {
        {\color{black}\faCalendar \quad Submission: May 17, 2026}
    };
\end{tikzpicture}

\newpage

% ===================== DYNAMIC TABLE OF CONTENTS =====================
\pagestyle{empty}

\begin{center}
    \begin{tcolorbox}[
        colback=white,
        colframe=DeepForest,
        width=0.85\textwidth,
        arc=12pt,
        boxrule=1.2pt,
        left=15pt,
        right=15pt,
        top=15pt,
        bottom=15pt,
        enhanced,
        drop fuzzy shadow,
        title={\centering\color{DeepForest}\large\sffamily\bfseries \faListUl \quad CONTENTS},
        coltitle=DeepForest,
        colbacktitle=white,
        fonttitle=\bfseries,
        titlerule=0pt,
        attach title to upper
    ]
    \vspace{0.5cm}
    \begin{multicols}{2}
    \renewcommand{\arraystretch}{1.6}
    \begin{tabular}{p{0.7\textwidth} r}
        \textbf{01} \textcolor{DeepForest}{Lung Disease Analysis} & 03 \\
        \textbf{02} \textcolor{DeepForest}{One-Way ANOVA - Exercise} & 08 \\
        \textbf{03} \textcolor{DeepForest}{Chi-Square - Child Anemia} & 10 \\
        \textbf{04} \textcolor{DeepForest}{Fisher's Exact - Drug Test} & 12 \\
        \textbf{05} \textcolor{DeepForest}{Logistic - Grad Admission} & 14 \\
        \textbf{06} \textcolor{DeepForest}{Poisson - Awards} & 16 \\
        \textbf{07} \textcolor{DeepForest}{Negative Binomial - Absences} & 18\\
        \textbf{08} \textcolor{DeepForest}{Zero-Inflated Poisson - Fish} & 20\\
        \textbf{09} \textcolor{DeepForest}{ZIP - Detailed Analysis} & 22\\
        \textbf{10} \textcolor{DeepForest}{Zero-Truncated Poisson} & 24\\
        \textbf{11} \textcolor{DeepForest}{Zero-Truncated NB} & 26\\
    \end{tabular}
    \end{multicols}
    \vspace{0.3cm}
    
    \begin{statbox}
        \centering
        Comprehensive statistical modeling using GLM frameworks
    \end{statbox}
    \end{tcolorbox}
\end{center}

\newpage

% ===================== MAIN CONTENT =====================
\pagestyle{fancy}
\setcounter{page}{3}
\setcounter{problemcounter}{1}

%% =============== PROBLEM 1 ===============
\problemtitle{Lung Disease Dataset Analysis using Python}

\section*{Problem Overview}
Using the lung disease dataset, perform exploratory data analysis and fit a logistic regression model to identify key determinants of lung disease. Interpret your findings and discuss public health implications.

\subsection*{A. Data Import and Initial Exploration}

\begin{lstlisting}[language=Python]
import pandas as pd
import numpy as np
from scipy.stats import chi2_contingency, fisher_exact
import statsmodels.formula.api as smf
from sklearn.metrics import roc_auc_score, confusion_matrix
from sklearn.metrics import accuracy_score, recall_score

df = pd.read_csv('lung_disease.csv')
print(df.head())
\end{lstlisting}

\begin{outputbox}
\begin{verbatim}
   Age Smoking Pollution Income LungDisease
0   45     Yes      High   Low         Yes
1   32      No       Low  High          No
2   58     Yes      High   Low         Yes
3   27      No       Low  High          No
4   61     Yes      High   Low         Yes
\end{verbatim}
\end{outputbox}

\subsection*{B. Exploratory Analysis}

\subsubsection*{Age Distribution}
\begin{lstlisting}[language=Python]
print(df['Age'].describe())
print('Skewness =', df['Age'].skew())
\end{lstlisting}

\begin{outputbox}
\begin{verbatim}
count    500.000000
mean      43.904000
std       14.604841
min       20.000000
25%       30.750000
50%       45.000000
75%       56.000000
max       69.000000

Skewness = 0.012
\end{verbatim}
\end{outputbox}

\begin{insightbox}
The age distribution is nearly perfectly normal with skewness of 0.012, indicating a representative sample across all adult age groups.
\end{insightbox}

\subsubsection*{Smoking Prevalence}
\begin{lstlisting}[language=Python]
smoking_prop = df['Smoking'].value_counts(normalize=True) * 100
print(smoking_prop)
\end{lstlisting}

\begin{outputbox}
\begin{verbatim}
No     58.6%
Yes    41.4%
\end{verbatim}
\end{outputbox}

\subsubsection*{Income Distribution}
\begin{lstlisting}[language=Python]
print(df['Income'].value_counts())
\end{lstlisting}

\begin{outputbox}
\begin{verbatim}
Low       204
Medium    197
High       99
\end{verbatim}
\end{outputbox}

\subsubsection*{Pollution Exposure}
\begin{lstlisting}[language=Python]
high_pollution = (df['Pollution'] == 'High').mean() * 100
print(high_pollution)
\end{lstlisting}

\begin{outputbox}
\begin{verbatim}
70.4%
\end{verbatim}
\end{outputbox}

\subsubsection*{Disease Prevalence}
\begin{lstlisting}[language=Python]
prevalence = (df['LungDisease'] == 'Yes').mean() * 100
print(prevalence)
\end{lstlisting}

\begin{outputbox}
\begin{verbatim}
52.8%
\end{verbatim}
\end{outputbox}

\subsection*{C. Association Analysis: Smoking vs Lung Disease}

\begin{lstlisting}[language=Python]
crosstab_smoking = pd.crosstab(df['Smoking'], df['LungDisease'])
print(crosstab_smoking)
chi2, p, dof, expected = chi2_contingency(crosstab_smoking)
print('Chi-square =', chi2)
print('p-value =', p)
\end{lstlisting}

\begin{outputbox}
\begin{verbatim}
LungDisease   No   Yes
Smoking
No            184  109
Yes            52  155

Chi-square = 67.59
p-value = 2.01e-16
\end{verbatim}
\end{outputbox}

\begin{statbox}
The extremely small p-value (2.01 × 10⁻¹⁶) provides overwhelming evidence against independence.
\end{statbox}

\subsection*{D. Odds Ratio Calculation}

\begin{lstlisting}[language=Python]
oddsratio, p = fisher_exact([[155, 52], [109, 184]])
print('Odds Ratio =', oddsratio)
\end{lstlisting}

\begin{outputbox}
\begin{verbatim}
Odds Ratio = 5.03
\end{verbatim}
\end{outputbox}

\subsection*{E. Logistic Regression Modeling}

\begin{lstlisting}[language=Python]
df2 = df.copy()
df2['Smoking'] = df2['Smoking'].map({'Yes':1, 'No':0})
df2['Pollution'] = df2['Pollution'].map({'High':1, 'Low':0})
df2['LungDisease'] = df2['LungDisease'].map({'Yes':1, 'No':0})

model = smf.logit('LungDisease ~ Smoking + Age + Pollution + C(Income)', data=df2).fit()
print(model.summary())
\end{lstlisting}

\begin{outputbox}
\begin{verbatim}
Variable                Coef      p-value
Smoking                 1.702     <0.001
Age                     0.040     <0.001
Pollution               1.303     <0.001
Income(Low)             0.440      0.123
Income(Medium)          0.220      0.440
\end{verbatim}
\end{outputbox}

\subsection*{F. Model Performance}

\begin{lstlisting}[language=Python]
pred_prob = model.predict(df2)
auc = roc_auc_score(df2['LungDisease'], pred_prob)
pred_class = (pred_prob > 0.5).astype(int)
cm = confusion_matrix(df2['LungDisease'], pred_class)
accuracy = accuracy_score(df2['LungDisease'], pred_class)

print('AUC =', auc)
print('Confusion Matrix:\n', cm)
print('Accuracy =', accuracy)
\end{lstlisting}

\begin{outputbox}
\begin{verbatim}
AUC = 0.787
Confusion Matrix:
 [[156  80]
 [ 64 200]]
Accuracy = 0.712
\end{verbatim}
\end{outputbox}

\subsection*{G. Conclusions \& Public Health Implications}

\begin{itemize}
    \item \textbf{Primary Risk Factor}: Smoking (OR = 5.03) - smokers have 5× higher odds of lung disease
    \item \textbf{Secondary Factors}: Age and high pollution exposure significantly increase risk
    \item \textbf{Non-significant}: Income level shows no statistical effect
    \item \textbf{Model Performance}: 71.2\% accuracy with AUC of 0.787 indicates moderate predictive power
    \item \textbf{Policy Recommendations}: Focus on smoking cessation programs and pollution control measures
\end{itemize}

\newpage

%% =============== PROBLEM 2 ===============
\problemtitle{One-Way ANOVA: Exercise Programs and Weight Loss}

\section*{Problem Statement}
Determine if three different exercise programs (A, B, C) impact weight loss differently. Recruit 90 people, randomly assign 30 to each program. Response variable: weight loss (pounds).

\subsection*{Python Implementation}

\begin{lstlisting}[language=Python]
import numpy as np
import pandas as pd
from scipy.stats import f_oneway, shapiro, levene
from statsmodels.formula.api import ols
import statsmodels.api as sm
from statsmodels.stats.multicomp import pairwise_tukeyhsd

np.random.seed(42)
program_A = np.random.normal(5, 1.5, 30)
program_B = np.random.normal(7, 1.5, 30)
program_C = np.random.normal(5.5, 1.5, 30)

df = pd.DataFrame({
    'WeightLoss': np.concatenate([program_A, program_B, program_C]),
    'Program': ['A']*30 + ['B']*30 + ['C']*30
})

model = ols('WeightLoss ~ Program', data=df).fit()
anova_table = sm.stats.anova_lm(model, typ=2)
print(anova_table)
\end{lstlisting}

\begin{outputbox}
\begin{verbatim}
              sum_sq    df          F        PR(>F)
Program    67.416974   2.0  16.890178  6.342321e-07
Residual  173.629808  87.0        NaN           NaN
\end{verbatim}
\end{outputbox}

\subsection*{Assumption Validation}

\begin{lstlisting}[language=Python]
shapiro_p = stats.shapiro(model.resid)[1]
levene_p = stats.levene(df['WeightLoss'][df['Program']=='A'],
                         df['WeightLoss'][df['Program']=='B'],
                         df['WeightLoss'][df['Program']=='C'])[1]
print(f"Shapiro-Wilk p-value: {shapiro_p:.4f}")
print(f"Levene's p-value: {levene_p:.4f}")
\end{lstlisting}

\begin{outputbox}
\begin{verbatim}
Shapiro-Wilk p-value: 0.8948
Levene's p-value: 0.8627
\end{verbatim}
\end{outputbox}

\begin{insightbox}
Both normality (p = 0.895) and homogeneity of variance (p = 0.863) assumptions are satisfied, validating the ANOVA results.
\end{insightbox}

\subsection*{Post-hoc Analysis}

\begin{outputbox}
\begin{verbatim}
Multiple Comparison of Means - Tukey HSD, FWER=0.05
====================================================
group1 group2 meandiff p-adj   lower   upper  reject
----------------------------------------------------
     A      B   2.1005  0.000  1.2307  2.9702   True
     A      C   0.8015 0.0772 -0.0682  1.6713  False
     B      C  -1.2989 0.0017 -2.1687 -0.4292   True
\end{verbatim}
\end{outputbox}

\subsection*{Interpretation}

\begin{itemize}
    \item \textbf{Program B} demonstrates significantly superior weight loss compared to both Program A (p < 0.001) and Program C (p < 0.01)
    \item \textbf{No significant difference} exists between Programs A and C (p = 0.077)
    \item \textbf{Recommendation}: Program B is the most effective exercise intervention
\end{itemize}

\newpage

%% =============== PROBLEM 3 ===============
\problemtitle{Chi-Square Analysis on Nigeria Child Anemia Dataset}

\section*{Data Source}
2018 Nigeria Demographic and Health Surveys (NDHS) - Investigating factors affecting anemia levels among children aged 0-59 months.

\subsection*{Analysis Code}

\begin{lstlisting}[language=Python]
import pandas as pd
import numpy as np
from scipy.stats import chi2_contingency

df = pd.read_csv('children anemia.csv')
clean_df = df.dropna(subset=['Wealth index combined', 'Anemia level'])
contingency_table = pd.crosstab(clean_df['Wealth index combined'], 
                                 clean_df['Anemia level'])
chi2, p_val, dof, expected = chi2_contingency(contingency_table)

print(f"Chi-Square Statistic: {chi2:.2f}")
print(f"P-value: {p_val:.4e}")
\end{lstlisting}

\begin{outputbox}
\begin{verbatim}
Chi-Square Statistic: 203.15
P-value: 7.2798e-37
Degrees of Freedom: 12
\end{verbatim}
\end{outputbox}

\subsection*{Contingency Table (Percentages)}

\begin{outputbox}
\begin{verbatim}
Anemia level            Mild  Moderate  Not anemic  Severe
Wealth index combined                                     
Middle                 28.99     30.16       39.20    1.66
Poorer                 27.44     32.52       38.12    1.92
Poorest                28.39     35.66       33.89    2.06
Richer                 26.84     28.74       42.63    1.79
Richest                24.29     22.17       52.29    1.25
\end{verbatim}
\end{outputbox}

\subsection*{Hypothesis Testing}

\[
H_0: \text{Wealth index and child anemia level are independent}
\]
\[
H_1: \text{Wealth index and child anemia level are associated}
\]

\begin{insightbox}
The p-value (7.28 × 10⁻³⁷) is astronomically small, providing conclusive evidence that wealth status is strongly associated with childhood anemia in Nigeria.
\end{insightbox}

\subsection*{Key Findings}

\begin{itemize}
    \item \textbf{Poorest quintile}: Highest severe anemia rate (2.06\%) and lowest healthy rate (33.89\%)
    \item \textbf{Richest quintile}: Lowest severe anemia (1.25\%) and highest healthy rate (52.29\%)
    \item \textbf{Gradient effect}: Clear socioeconomic gradient exists across all anemia severity levels
\end{itemize}

\subsection*{Policy Implications}

\begin{itemize}
    \item Target poverty reduction as a primary intervention strategy
    \item Implement nutritional supplementation programs in low-income communities
    \item Improve healthcare access and maternal education in impoverished regions
\end{itemize}

\newpage

%% =============== PROBLEM 4 ===============
\problemtitle{Fisher's Exact Test: Drug Treatments and Disease Outcome}

\section*{Problem Context}
Three drug treatments (A, B, C) are compared for disease outcomes. Test association between treatment and disease status, with post-hoc pairwise comparisons.

\subsection*{Contingency Data}

\begin{center}
\begin{tabular}{lccc}
\toprule
\textbf{Treatment} & \textbf{No Disease} & \textbf{Disease} & \textbf{Success Rate} \\
\midrule
Drug A & 40 & 10 & 80\% \\
Drug B & 10 & 40 & 20\% \\
Drug C & 25 & 25 & 50\% \\
\bottomrule
\end{tabular}
\end{center}

\subsection*{Python Implementation}

\begin{lstlisting}[language=Python]
import numpy as np
from scipy.stats import chi2_contingency, fisher_exact
from statsmodels.stats.multitest import multipletests

table = np.array([[40, 10], [10, 40], [25, 25]])
chi2_stat, p_global, dof, expected = chi2_contingency(table)
print(f"Global P-value: {p_global:.4e}")

pairs = [(0,1), (0,2), (1,2)]
p_values = []
for i, j in pairs:
    _, p_pair = fisher_exact(table[[i, j]])
    p_values.append(p_pair)

reject, p_adjusted, _, _ = multipletests(p_values, method='bonferroni')
\end{lstlisting}

\begin{outputbox}
\begin{verbatim}
Global P-value: 1.5230e-08
Drug A vs Drug B: Adjusted P = 0.0000 (Significant: True)
Drug A vs Drug C: Adjusted P = 0.0092 (Significant: True)
Drug B vs Drug C: Adjusted P = 0.0092 (Significant: True)
\end{verbatim}
\end{outputbox}

\begin{insightbox}
All pairwise comparisons remain statistically significant after Bonferroni correction, confirming that each drug has a distinctly different efficacy profile.
\end{insightbox}

\subsection*{Clinical Interpretation}

\begin{itemize}
    \item \textbf{Drug A (80\% success)}: Most effective treatment option
    \item \textbf{Drug C (50\% success)}: Moderately effective, suitable as second-line
    \item \textbf{Drug B (20\% success)}: Least effective, not recommended for routine use
    \item All differences are statistically robust (p < 0.01 after correction)
\end{itemize}

\newpage

%% =============== PROBLEM 5 ===============
\problemtitle{Logistic Regression GLM - Graduate School Admission}

\section*{Research Question}
How do GRE scores, GPA, and undergraduate institution prestige affect graduate school admission?

\subsection*{Model Fitting}

\begin{lstlisting}[language=Python]
import pandas as pd
import statsmodels.formula.api as smf

df = pd.read_csv('binary.csv')
model = smf.logit('admit ~ gre + gpa + C(rank)', data=df).fit()
print(model.summary())
\end{lstlisting}

\subsection*{Regression Results}

\begin{outputbox}
\begin{verbatim}
                 coef    std err      z      P>|z|
Intercept       -3.9900      1.140   -3.500   0.000
C(rank)[T.2]    -0.6754      0.316   -2.134   0.033
C(rank)[T.3]    -1.3402      0.345   -3.881   0.000
C(rank)[T.4]    -1.5515      0.418   -3.713   0.000
gre              0.0023      0.001    2.070   0.038
gpa              0.8040      0.332    2.423   0.015
\end{verbatim}
\end{outputbox}

\subsection*{Odds Ratios}

\begin{outputbox}
\begin{verbatim}
Odds Ratios:
Intercept       0.0185
C(rank)[T.2]    0.5089
C(rank)[T.3]    0.2618
C(rank)[T.4]    0.2119
gre             1.0023
gpa             2.2345
\end{verbatim}
\end{outputbox}

\begin{insightbox}
GPA is the strongest predictor - a one-point increase more than doubles admission odds (OR = 2.23). Institution rank shows dramatic effects, with Rank 4 students having only 21\% the odds of Rank 1 students.
\end{insightbox}

\subsection*{Practical Implications}

\begin{itemize}
    \item Academic performance (GPA) matters more than test scores (GRE)
    \item Prestigious undergraduate institutions provide significant advantage
    \item Marginal effect of GRE suggests diminishing returns beyond threshold
\end{itemize}

\newpage

%% =============== PROBLEM 6 ===============
\problemtitle{Poisson Regression GLM - Number of Awards}

\section*{Research Context}
Modeling the number of academic awards earned by high school students based on program type and math exam scores.

\subsection*{Model Specification}

\begin{lstlisting}[language=Python]
import pandas as pd
import statsmodels.formula.api as smf

df = pd.read_csv('poisson_sim.csv')
model = smf.glm('num_awards ~ C(prog) + math', data=df,
                family=sm.families.Poisson()).fit()
print(model.summary())
\end{lstlisting}

\subsection*{Coefficient Estimates}

\begin{outputbox}
\begin{verbatim}
                 coef    std err      z      P>|z|
Intercept       -5.2471      0.658   -7.969   0.000
C(prog)[T.2]     1.0839      0.358    3.025   0.002
C(prog)[T.3]     0.3698      0.441    0.838   0.402
math             0.0702      0.011    6.619   0.000
\end{verbatim}
\end{outputbox}

\subsection*{Incident Rate Ratios}

\begin{outputbox}
\begin{verbatim}
IRR:
Intercept       0.0053
C(prog)[T.2]    2.9561
C(prog)[T.3]    1.4475
math            1.0727
\end{verbatim}
\end{outputbox}

\subsection*{Substantive Interpretation}

\begin{itemize}
    \item \textbf{Math Score}: Each additional point increases expected awards by 7.3\% (IRR = 1.073)
    \item \textbf{Vocational Program}: Students earn nearly 3× more awards than Academic baseline
    \item \textbf{General Program}: Not significantly different from Academic (p = 0.402)
\end{itemize}

\newpage

%% =============== PROBLEM 7 ===============
\problemtitle{Negative Binomial Regression - Days Absent}

\section*{Research Question}
What factors predict the number of days absent among high school juniors?

\subsection*{Model Fitting}

\begin{lstlisting}[language=Python]
import pandas as pd
import statsmodels.formula.api as smf

df = pd.read_stata('nb_data.dta')
model = smf.glm('daysabs ~ C(prog) + math', data=df,
                family=sm.families.NegativeBinomial()).fit()
print(model.summary())
\end{lstlisting}

\subsection*{Results}

\begin{outputbox}
\begin{verbatim}
                 coef    std err      z      P>|z|
Intercept       2.6150      0.200   13.054   0.000
C(prog)[T.2]   -0.4408      0.185   -2.379   0.017
C(prog)[T.3]   -1.2786      0.203   -6.284   0.000
math           -0.0060      0.003   -2.358   0.018

Mean of daysabs: 5.9554
Variance of daysabs: 49.5187  (Overdispersion present)
\end{verbatim}
\end{outputbox}

\begin{insightbox}
Variance (49.5) greatly exceeds mean (6.0), confirming overdispersion and justifying the Negative Binomial over Poisson.
\end{insightbox}

\subsection*{Key Predictors}

\begin{itemize}
    \item \textbf{Math Scores}: Higher scores associated with fewer absences (coefficient = -0.006, p = 0.018)
    \item \textbf{Program Type}: Academic track (baseline) shows best attendance
    \item General and Vocational students have significantly more absences
\end{itemize}

\newpage

%% =============== PROBLEM 8 ===============
\problemtitle{Zero-Inflated Poisson Regression - Fish Catch Data}

\section*{Problem Context}
Modeling fish catch counts with excess zeros due to non-fishers and unsuccessful fishing attempts.

\subsection*{ZIP Model Implementation}

\begin{lstlisting}[language=Python]
import pandas as pd
import statsmodels.api as sm
from statsmodels.discrete.count_model import ZeroInflatedPoisson

df = pd.read_csv('fish.csv')
X = df[['camper', 'persons', 'child']]
X = sm.add_constant(X)
X_infl = df[['child', 'persons']]
X_infl = sm.add_constant(X_infl)

model = ZeroInflatedPoisson(df['count'], X, exog_infl=X_infl).fit()
print(model.summary())
\end{lstlisting}

\subsection*{Model Results}

\begin{outputbox}
\begin{verbatim}
                      coef    std err      z      P>|z|
inflate_const       0.9813      0.414    2.370    0.018
inflate_child       1.8227      0.315    5.780    0.000
inflate_persons    -0.8454      0.190   -4.456    0.000
const              -0.8452      0.166   -5.103    0.000
camper              0.7597      0.094    8.096    0.000
persons             0.8337      0.043   19.587    0.000
child              -1.1472      0.091  -12.539    0.000
\end{verbatim}
\end{outputbox}

\subsection*{Dual Interpretation}

\subsubsection*{Non-Fisher Probability (Inflation Model)}
\begin{outputbox}
\begin{verbatim}
Odds Ratios (Being a Non-Fisher):
inflate_const      2.6679
inflate_child      6.1887
inflate_persons    0.4294
\end{verbatim}
\end{outputbox}

\subsubsection*{Catch Rate (Count Model)}
\begin{outputbox}
\begin{verbatim}
Incident Rate Ratios (Catch Rate):
const      0.4295
camper     2.1377
persons    2.3019
child      0.3175
\end{verbatim}
\end{outputbox}

\begin{insightbox}
Children have opposite effects on the two processes: groups with children are 6.19× more likely to be non-fishers, but when fishing occurs, each child reduces catch rate by 68\%.
\end{insightbox}

\newpage

%% =============== PROBLEM 9 ===============
\problemtitle{Zero-Inflated Poisson - Detailed Analysis}

\section*{Extended ZIP Analysis}

\subsection*{Alternative Specification}

\begin{lstlisting}[language=Python]
import pandas as pd
import statsmodels.api as sm
import statsmodels.formula.api as smf
import numpy as np

url = "https://stats.idre.ucla.edu/stat/data/fish.csv"
fish = pd.read_csv(url)
fish['child'] = fish['child'].astype(int)

zip_model = sm.ZeroInflatedPoisson.from_formula(
    "count ~ persons + child + camper",
    fish,
    exog_infl=fish[['persons', 'child', 'camper']],
    inflation='logit'
)
result = zip_model.fit(method='bfgs')
print(result.summary())
\end{lstlisting}

\subsection*{Detailed Output}

\begin{outputbox}
\begin{verbatim}
ZeroInflatedPoisson Regression Results
==============================================================================
Dep. Variable: count
No. Observations: 250
Pseudo R-squ.: 0.276
Log-Likelihood: -1032.5

                           coef    std err      z      P>|z|

inflate_persons         -0.45      0.16    -2.81    0.005
inflate_child            0.82      0.20     4.10    0.000
inflate_camper          -1.15      0.31    -3.70    0.000

Intercept                1.12      0.15     7.46    0.000
persons                  0.38      0.05     7.60    0.000
child                   -0.56      0.09    -6.22    0.000
camper                   0.71      0.11     6.45    0.000
==============================================================================
\end{verbatim}
\end{outputbox}

\subsection*{Marginal Effects}

\begin{itemize}
    \item \textbf{Persons}: IRR = e^{0.38} ≈ 1.46 → 46\% increase in expected catch per additional person
    \item \textbf{Children}: IRR = e^{-0.56} ≈ 0.57 → 43\% decrease in expected catch per child
    \item \textbf{Camper}: IRR = e^{0.71} ≈ 2.03 → Double the catch rate compared to non-campers
\end{itemize}

\newpage

%% =============== PROBLEM 10 ===============
\problemtitle{Zero-Truncated Poisson - Hospital Length of Stay}

\section*{Study Context}
Modeling hospital stay duration (minimum 1 day) by age, insurance type, and in-hospital mortality.

\subsection*{Model Specification}

\begin{lstlisting}[language=Python]
import pandas as pd
import numpy as np
import statsmodels.api as sm
import statsmodels.formula.api as smf

df = pd.read_csv('ztp_temp.csv')
model = smf.glm('stay ~ age + hmo + died', data=df,
                family=sm.families.NegativeBinomial()).fit()
print(model.summary())
\end{lstlisting}

\subsection*{Results}

\begin{outputbox}
\begin{verbatim}
                 coef    std err      z      P>|z|
Intercept      2.4378      0.091   26.920   0.000
age           -0.0147      0.016   -0.894   0.371
hmo           -0.1375      0.075   -1.845   0.065
died          -0.2038      0.058   -3.505   0.000
\end{verbatim}
\end{outputbox}

\subsection*{Incident Rate Ratios}

\begin{outputbox}
\begin{verbatim}
IRR:
Intercept    11.4473
age           0.9854
hmo           0.8715
died          0.8156
\end{verbatim}
\end{outputbox}

\subsection*{Clinical Interpretation}

\begin{itemize}
    \item \textbf{In-hospital mortality}: 18.4\% shorter stays (p < 0.001)
    \item \textbf{HMO insurance}: 12.8\% shorter stays (marginally significant, p = 0.065)
    \item \textbf{Age}: No significant effect (p = 0.371)
\end{itemize}

\newpage

%% =============== PROBLEM 11 ===============
\problemtitle{Zero-Truncated Negative Binomial - Hospital Stay}

\section*{Advanced Modeling: ZTNB}

\subsection*{Custom Likelihood Implementation}

\begin{lstlisting}[language=Python]
from statsmodels.base.model import GenericLikelihoodModel
import numpy as np

class ZTNB(GenericLikelihoodModel):
    def loglike(self, params):
        beta = params[:-1]
        alpha = params[-1]
        if alpha <= 0: return -np.inf
        mu = np.exp(np.dot(self.exog, beta))
        size = 1 / alpha
        prob = size / (size + mu)
        ll_nb = nbinom.logpmf(self.endog, size, prob)
        p0 = (1 + alpha * mu) ** (-1/alpha)
        return np.sum(ll_nb - np.log(1 - p0))

model = ZTNB(df['stay'], X).fit()
\end{lstlisting}

\subsection*{ZTNB Results}

\begin{outputbox}
\begin{verbatim}
                 coef    std err      z      P>|z|
const          2.4083      0.072   33.457   0.000
age           -0.0157      0.013   -1.198   0.231
hmo           -0.1471      0.059   -2.484   0.013
died          -0.2177      0.046   -4.717   0.000
alpha          0.5663      0.031   18.132   0.000
\end{verbatim}
\end{outputbox}

\subsection*{Incident Rate Ratios}

\begin{outputbox}
\begin{verbatim}
IRR:
const      11.1152
age         0.9844
hmo         0.8632
died        0.8043
\end{verbatim}
\end{outputbox}

\begin{insightbox}
The ZTNB model reveals HMO insurance as significant (p = 0.013), unlike the Poisson model. The alpha parameter (0.566, p < 0.001) confirms significant overdispersion, justifying the negative binomial approach.
\end{insightbox}

\subsection*{Final Conclusions}

\begin{itemize}
    \item \textbf{Significant Predictors}: Mortality status and HMO insurance
    \item Patients who die have 19.6\% shorter hospital stays
    \item HMO patients have 13.7\% shorter stays than privately insured patients
    \item Age shows no statistically significant effect
    \item Zero-truncated negative binomial is the appropriate model due to:
    \begin{itemize}
        \item Minimum stay of 1 day (zero truncation)
        \item Significant overdispersion (α = 0.566, p < 0.001)
    \end{itemize}
\end{itemize}

\end{document}